\section{Cálculo discreto de la medida sobre el límite inverso: diferencias, cociclo de acarreo, ultramétrica, reconstrucción de \texorpdfstring{\(\mathbb R\)}{R} y desintegración genealógica}
\label{sec:complejo-discreto-medida-rev2}
\index{medida!complejo discreto}
\index{camino!integración}
\index{acarreo!cociclo}

La distinción entre contar y medir admite una formulación operatoria. Contar
asigna una cardinalidad a una colección de marcas. Medir evalúa relaciones
entre marcas mediante una regla local, las integra a lo largo de un camino y
conserva la carta en la que se publica el resultado. La unidad mínima de ese
proceso es una diferencia; su composición produce longitud, circulación y,
cuando interviene una reducción modular, acarreo. Esta arquitectura sitúa la
medida antes de su coordenada real: el escalar publicado es la contracción de
un cálculo que posee dominio, orientación y memoria.

\subsection{Marcas y operador de diferencias}

Sea \(P_n\) el camino con vértices ordenados
\(V(P_n)=\{v_0,\ldots,v_{n-1}\}\) y aristas orientadas
\(e_j=(v_{j-1},v_j)\), \(1\leq j\leq n-1\). Para un anillo conmutativo
\(A\), definimos
\[
 C^0(P_n;A)=A^{V(P_n)},
 \qquad
 C^1(P_n;A)=A^{E(P_n)},
\]
y el operador de diferencias
\begin{equation}
 d:C^0(P_n;A)\longrightarrow C^1(P_n;A),
 \qquad
 (df)(e_j)=f(v_j)-f(v_{j-1}).
 \label{eq:diferencia-cero-uno-rev2}
\end{equation}

\begin{theorem}[Rango del operador de diferencias sobre un camino]
\label{thm:rango-diferencias-camino-rev2}
Si \(A\) es un cuerpo, entonces
\[
 \ker d=A\mathbf 1,
 \qquad
 \operatorname{rank}d=n-1.
\]
Por consiguiente, \(n\) marcas conectadas determinan \(n-1\) diferencias
independientes y un único modo constante.
\end{theorem}

\begin{proof}
La igualdad \(df=0\) implica sucesivamente
\(f(v_0)=f(v_1)=\cdots=f(v_{n-1})\), de modo que el núcleo está formado por
las funciones constantes y tiene dimensión uno. El teorema de
rango--nulidad da \(\operatorname{rank}d=n-1\). Equivalentemente, toda
asignación \((a_1,\ldots,a_{n-1})\) a las aristas se integra de manera única
una vez fijado \(f(v_0)\):
\[
 f(v_k)=f(v_0)+\sum_{j=1}^{k}a_j.
\]
\end{proof}

El modo constante expresa la libertad de elegir origen. Las diferencias
contienen la información relacional; la marca inicial proporciona la sección
que reconstruye las posiciones. En la célula inaugural de HMT, nueve marcas
ordenadas producen ocho diferencias interiores. El cierre periódico añade la
novena arista y transforma el camino en ciclo; la completación bidimensional
de esos intervalos se realizará de forma exacta en la construcción de APP.

\subsection{Longitud positiva e integral orientada}

Dos evaluaciones de una ruta deben permanecer separadas. Una métrica discreta
es una cocadena positiva
\[
 \ell:E(P_n)\longrightarrow\mathbb R_{>0}.
\]
Para un camino no orientado \(|\gamma|\), su longitud es
\begin{equation}
 L_\ell(\gamma)=\sum_{e\in|\gamma|}m_\gamma(e)\,\ell(e),
 \label{eq:longitud-positiva-camino-rev2}
\end{equation}
donde \(m_\gamma(e)\) cuenta las travesías de la arista. Una cocadena
orientada \(\omega\in C^1(P_n;A)\), en cambio, se integra con signo:
\begin{equation}
 \int_\gamma\omega
 :=\sum_{k=1}^{r}\varepsilon_k\omega(e_{i_k}),
 \qquad \varepsilon_k\in\{-1,+1\}.
 \label{eq:integral-cocadena-camino-rev2}
\end{equation}
La longitud suma coste geométrico; la integral orientada suma transporte. La
primera es invariante bajo inversión de la ruta y la segunda cambia de signo.

\begin{proposition}[Teorema fundamental discreto]
\label{prop:teorema-fundamental-discreto-rev2}
Para \(f\in C^0(P_n;A)\) y toda ruta orientada \(\gamma:a\to b\),
\[
 \int_\gamma df=f(b)-f(a).
\]
En particular, la integral de una cocadena exacta sobre un ciclo es nula.
\end{proposition}

\begin{proof}
La suma de las diferencias consecutivas es telescópica. Cada vértice interior
aparece una vez con signo positivo y una vez con signo negativo; permanecen
únicamente los extremos.
\end{proof}

La proposición distingue dos niveles de cierre. El grupo
\(H_1(C_n;\mathbb Z)\cong\mathbb Z\) proporciona el portador topológico de la
vuelta; una holonomía requiere además un transporte. Si
\(T:\operatorname{Path}(C_n)\to\mathcal C\) es un funtor hacia una categoría
de automorfismos, la holonomía de un ciclo \(\gamma\) es
\begin{equation}
 \operatorname{Hol}_T(\gamma)=T(e_r)\cdots T(e_1).
 \label{eq:holonomia-transporte-ciclo-rev2}
\end{equation}
El ciclo aporta la clase de recorrido; \(T\) aporta la conexión discreta que
convierte esa clase en un endomorfismo. Esta distinción será esencial cuando
el retorno visible recupere la fase y el transporte conserve un estado de
frontera diferente.

\subsection{Acarreo como cociclo de composición}

La reducción modular publica una clase y una sección elige representantes.
Sea \(Q=\mathbb Z/9\mathbb Z\) y sea
\(s_0:Q\to\{0,\ldots,8\}\subset\mathbb Z\) la sección residual ordinaria.
Definimos
\begin{equation}
 \kappa_9(a,b)
 :=\frac{s_0(a)+s_0(b)-s_0(a+b)}{9}\in\mathbb Z.
 \label{eq:cociclo-acarreo-nonadico-rev2}
\end{equation}
Entonces
\[
 s_0(a)+s_0(b)=s_0(a+b)+9\kappa_9(a,b).
\]

\begin{proposition}[Identidad cocíclica del acarreo]
\label{prop:acarreo-dos-cociclo-rev2}
La aplicación \(\kappa_9:Q\times Q\to\mathbb Z\) satisface
\begin{equation}
 \kappa_9(a,b)+\kappa_9(a+b,c)
 =\kappa_9(b,c)+\kappa_9(a,b+c).
 \label{eq:identidad-dos-cociclo-rev2}
\end{equation}
Es, por tanto, un \(2\)-cociclo normalizado asociado a la sección elegida.
\end{proposition}

\begin{proof}
Se expanden ambos miembros mediante
\eqref{eq:cociclo-acarreo-nonadico-rev2}. Tras multiplicar por nueve, los dos
son iguales a
\[
 s_0(a)+s_0(b)+s_0(c)-s_0(a+b+c).
\]
La normalización se obtiene de \(s_0(0)=0\).
\end{proof}

La identidad cocíclica expresa la asociatividad del levantamiento: agrupar primero
\((a+b)+c\) o \(a+(b+c)\) produce la misma clase visible y el mismo acarreo
total. En una categoría de rutas, la misma estructura adopta la forma
\[
 \widetilde g\,\widetilde h
 =\kappa(g,h)\,\widetilde{gh},
\]
donde \(\widetilde g,\widetilde h\) son transportes enriquecidos que levantan
morfismos visibles. El acarreo es así un defecto de composición perfectamente
controlado antes de recibir cualquier interpretación geométrica como
curvatura o torsión.

\subsection{Compatibilidad de las cartas nonádica, ternaria y de completación decimal en base \texorpdfstring{\(10^3\)}{10^3}}
\label{subsec:compatibilidad-lectores-9-3-1000-rev3}

Los lectores residuales que intervienen en TPK no son nombres distintos de
una misma operación. Sea
\[
 R_9:=\mathbb Z/9\mathbb Z,
 \qquad
 r_9:\mathbb Z\longrightarrow R_9,
 \qquad
 r_3:\mathbb Z\longrightarrow\mathbb F_3
\]
las reducciones canónicas, y sea
\[
 \pi_3:R_9\longrightarrow\mathbb F_3,
 \qquad \pi_3(\overline n)=n\bmod3.
\]
El cuadrado de reducciones conmuta:
\begin{equation}
\begin{CD}
 \mathbb Z @>{r_9}>> R_9\\
 @V{r_3}VV @VV{\pi_3}V\\
 \mathbb F_3 @= \mathbb F_3 .
\end{CD}
\qquad
 r_3=\pi_3\circ r_9.
\label{eq:diagrama-reducciones-9-3-rev3}
\end{equation}
La sección positiva
\[
 s_+:R_9\longrightarrow\{1,\ldots,9\},
 \qquad
 s_+(\overline0)=9,
 \qquad
 s_+(\overline k)=k\quad(1\le k\le8),
\]
produce, para \(n>0\), la raíz digital positiva
\[
 \operatorname{dr}_9^+(n)=s_+(r_9(n)).
\]
Esta sección no es un homomorfismo: publica la clase nula mediante \(9\) y
olvida el cociente. La sección residual \(s_0(\overline0)=0\) publica esa
misma clase mediante \(0\). Por tanto, el cambio visible \(9\mapsto0\) es un
cambio de sección, y la igualdad
\[
 9=s_0(\overline0)+9\cdot1
\]
conserva como cociente el acarreo \(1\). El sucesor topológico del ciclo
positivo es, en cambio, \(9\mapsto1\), con
\[
 10=s_+(\overline1)+9\cdot1.
\]
Así se separan el cierre \(P_9\to C_9\), el cambio de carta y el incremento
de memoria; ninguno de los tres se sustituye por los otros dos.

La coordenada del radical pertenece a otro diagrama. Para
\(\mathfrak m=(3)\subset R_9\), la aplicación
\[
 \iota_{\mathfrak m}:\mathbb F_3\longrightarrow\mathfrak m,
 \qquad \overline a\longmapsto3a,
\]
es un isomorfismo de espacios vectoriales sobre \(\mathbb F_3\), con inversa
\begin{equation}
 \eta_{\mathfrak m}:\mathfrak m\longrightarrow\mathbb F_3,
 \qquad
 \eta_{\mathfrak m}(3a)=a\bmod3.
 \label{eq:coordenada-radical-compatibilidad-rev3}
\end{equation}
No es un isomorfismo de anillos. Tampoco coincide con la reducción directa:
para \(3a\in\mathfrak m\), se tiene \(\pi_3(3a)=0\), mientras
\(\eta_{\mathfrak m}(3a)=a\bmod3\). La primera lectura publica la clase
residual; la segunda publica la coordenada interna del radical.

Finalmente, sean \(u_j\in\{0,\ldots,999\}\) y
\[
 D_N=1000D_{N-1}+u_N,
 \qquad D_0=0,
\]
la concatenación de bloques en base \(10^3\). Como \(10^3\equiv1\pmod9\),
\begin{equation}
 r_9(D_N)=r_9\!\left(\sum_{j=1}^{N}u_j\right),
 \qquad
 r_3(D_N)=r_3\!\left(\sum_{j=1}^{N}u_j\right).
 \label{eq:compatibilidad-base-mil-reducciones-rev3}
\end{equation}
La base \(10^3\) conserva el orden de los bloques; las reducciones conservan sólo
su clase acumulada. La raíz digital añade una sección positiva; la coordenada
radical recupera una dirección interna de \((3)\). Esta jerarquía fija qué
información publica cada lector y qué cociente, orden o carta debe retener la
memoria TPK para reconstruir el antecedente.

\subsection{Sistema graduado de medida y registro no contraído}

Los tres niveles anteriores se organizan en un sistema graduado:
\[
 \boxed{
 \begin{array}{ccl}
 \text{grado }0&:&\text{valores o estados sobre marcas},\\
 \text{grado }1&:&\text{diferencias y transportes sobre aristas},\\
 \text{grado }2&:&\text{cociclos de composición y acarreo}.
 \end{array}}
\]
Los grados cero y uno forman el complejo celular
\(C^0\xrightarrow d C^1\); el grado dos vive en el complejo del nervio de la
categoría de caminos. Su reunión constituye el \emph{complejo discreto de
medida}: una estructura organizadora que conserva la diferencia entre
geometría celular y memoria categorial.

Para una ruta enriquecida
\(\Gamma=(x_0\xrightarrow{e_1}x_1\to\cdots\xrightarrow{e_r}x_r)\), el
registro es la integral no contraída
\begin{equation}
 \Lambda(\Gamma)
 =\bigl((x_{k-1},e_k,x_k),\omega_k,\kappa_k,\sigma_k,b_k,m_k\bigr)_{k=1}^{r},
 \label{eq:ledger-integral-no-contraida-rev2}
\end{equation}
donde \(\omega_k\) es la lectura local, \(\kappa_k\) el acarreo, \(\sigma_k\)
la orientación o firma parcial, \(b_k\) el dato de frontera y \(m_k\) la
memoria conservada. Un observable es una aplicación tipada
\[
 \mathcal O=F_{\mathcal O}\circ\Lambda.
\]
La evaluación puede publicar una suma, una clase residual, una firma o un
escalar sin convertir esas salidas en estados completos. La concatenación de
rutas concatena registros y añade el cociclo en la frontera común; por ello el
registro no contraído conserva exactamente la información necesaria para recomponer la
historia que una coordenada aislada contrae.

\begin{figure}[tbp]
\centering
\begin{tikzpicture}[
  >=Latex,
  every node/.style={font=\small,align=center},
  object/.style={draw,rounded corners=2pt,inner sep=5pt,text width=3.00cm,
                 minimum height=1.25cm},
  arrow/.style={->,semithick},
  maplabel/.style={font=\footnotesize,fill=white,inner sep=1.5pt}
]
\node[object] (marks) at (0,0)
  {marcas y estados\\\(C^0\)};
\node[object] (edges) at (5.25,0)
  {diferencias y transporte\\\(C^1\)};
\node[object] (paths) at (10.50,0)
  {caminos y ciclos\\integración y holonomía};
\node[object] (memory) at (0,-2.35)
  {cociclo de composición\\acarreo y frontera};
\node[object] (section) at (5.25,-2.35)
  {registro no contraído\\sección compatible};
\node[object] (outputs) at (10.50,-2.35)
  {evaluaciones tipadas\\valor, firma o simetría};
\draw[arrow] (marks) -- node[maplabel,above]{\(d\)} (edges);
\draw[arrow] (edges) -- node[maplabel,above]{\(\int_\gamma\)} (paths);
\draw[arrow] (paths) -- node[maplabel,right]{transporte TPK} (outputs);
\draw[arrow] (memory) -- node[maplabel,above]{composición} (section);
\draw[arrow] (section) -- node[maplabel,above]{contracción} (outputs);
\draw[arrow] (edges) -- node[maplabel,left]{defecto de levantamiento} (memory);
\draw[arrow] (paths) -- node[maplabel,right]{historia} (section);
\end{tikzpicture}
\caption{Arquitectura del complejo discreto de medida. El operador de
diferencias, la integral de caminos y el cociclo de composición ocupan
dominios distintos; el TPK los reúne en una sección compatible antes de que
un lector publique un valor o una simetría.}
\label{fig:complejo-discreto-medida-rev3}
\end{figure}

\subsection{Reconstrucción del continuo a partir del sistema inverso de cilindros compatibles}

La prolongación a profundidad arbitraria se formula mediante un sistema
inverso. Sean \(S_n\) conjuntos finitos no vacíos de prefijos enriquecidos y
\(r_{n+1,n}:S_{n+1}\to S_n\) aplicaciones de restricción sobreyectivas. Su
espacio de secciones compatibles es
\[
 \Omega_\infty=\varprojlim_n S_n
 =\{(x_n)_n:r_{n+1,n}(x_{n+1})=x_n\}.
\]
Cada \(S_n\) se dota de la topología discreta y \(\Omega_\infty\) de la
topología de subespacio del producto.

\subsection{Cilindros, ultramétrica y retorno arquimediano con conservación de memoria}

Para \(a\in S_n\), el cilindro de prefijo \(a\) es
\begin{equation}
 [a]_n:=\{x\in\Omega_\infty:x_n=a\}.
 \label{eq:cilindro-genealogico-rev3}
\end{equation}
Dos cilindros del mismo nivel son disjuntos, y cada cilindro de nivel
\(n+1\) queda contenido en el cilindro de su restricción. Esta base de
abiertos y cerrados publica la profundidad compartida por dos secciones sin
contraer su memoria a una coordenada real.

Para \(x\neq y\), sea
\[
 N(x,y):=\min\{n\geq0:x_n\neq y_n\}.
\]
Fijado \(0<\vartheta<1\), definimos
\begin{equation}
 d_\vartheta(x,y):=\vartheta^{N(x,y)},
 \qquad d_\vartheta(x,x):=0.
 \label{eq:ultrametrica-genealogica-rev3}
\end{equation}

\begin{proposition}[Ultramétrica de primera divergencia]
\label{prop:ultrametrica-genealogica-rev3}
La función \(d_\vartheta\) es una ultramétrica y sus bolas son precisamente
los cilindros, con el reajuste de índice correspondiente al nivel inicial:
\[
 d_\vartheta(x,z)
 \leq
 \max\{d_\vartheta(x,y),d_\vartheta(y,z)\}.
\]
Por tanto, dos secciones que publican el mismo valor pueden permanecer a
distancia positiva cuando divergen su hoja, su acarreo, su frontera o su
historia de prolongación.
\end{proposition}

\begin{proof}
Si \(x\) e \(y\) coinciden hasta el nivel \(n\), y \(y\) y \(z\) también,
entonces \(x\) y \(z\) coinciden hasta ese nivel. La primera divergencia del
primer y tercer estado ocurre, por consiguiente, no antes que el mínimo de
las otras dos. Como \(0<\vartheta<1\), elevar \(\vartheta\) a esos
exponentes produce la desigualdad fuerte. La descripción de las bolas sigue
directamente de fijar un prefijo común.
\end{proof}

\begin{theorem}[Compacidad y completitud del espacio genealógico]
\label{thm:compacidad-ultrametrica-genealogica-rev3}
Si cada \(S_n\) es finito y los enlaces son sobreyectivos, entonces
\(\Omega_\infty\), con \(d_\vartheta\), es no vacío, compacto y completo.
Si además
\begin{equation}
 \forall x\in\Omega_\infty\ \forall n\ \exists y\neq x:
 y_n=x_n,
 \label{eq:perfeccion-genealogica-rev3}
\end{equation}
no posee puntos aislados. En ese régimen es un compacto metrizable,
cero-dimensional y perfecto; por la caracterización clásica, es homeomorfo
al espacio de Cantor.
\end{theorem}

\begin{proof}
La sobreyectividad y la ramificación finita proporcionan una sección por el
lema de König. El producto de los conjuntos finitos discretos es compacto, y
las ecuaciones \(r_{n+1,n}(x_{n+1})=x_n\) definen un subespacio cerrado. La
topología inducida coincide con la topología de cilindros y, por la
\cref{prop:ultrametrica-genealogica-rev3}, con la de
\(d_\vartheta\). Todo espacio métrico compacto es completo. La condición
\eqref{eq:perfeccion-genealogica-rev3} garantiza que cada cilindro contiene
una sección distinta de su centro, por lo que ningún singleton es abierto.
\end{proof}

La periodicidad visible y el retorno del estado pertenecen a mapas
distintos. Sea \(p:E\twoheadrightarrow B\) una proyección, \(T:E\to E\) un
transporte enriquecido y \(f:B\to B\) la evolución publicada, con
\(p\circ T=f\circ p\). Si \(f^9=I_B\), entonces
\begin{equation}
 p\circ T^9=p,
 \label{eq:retorno-visible-fibra-rev3}
\end{equation}
de modo que \(T^9\) preserva cada fibra de \(p\), pero no tiene por qué ser
la identidad en ella. En la coordenada nonádica
\[
 (q,r)\longmapsto
 \begin{cases}
  (q,r+1),&r<9,\\
  (q+1,1),&r=9,
 \end{cases}
\]
una vuelta recupera la posición y aumenta la altura. En la periodicidad
CT108, nueve iteraciones avanzan una fase dodecafásica y ciento ocho
recuperan posición y fase; el estado enriquecido puede conservar todavía
hoja, frontera, orientación y genealogía diferentes. Ésta es la distinción
formal entre retorno de fase y retorno de estado.

\begin{theorem}[Reconstrucción del valor por secciones compatibles]
\label{thm:reconstruccion-continuo-secciones-rev2}
El espacio \(\Omega_\infty\) es no vacío y compacto. Supóngase que a cada
prefijo \(x_n\in S_n\) se asigna un intervalo compacto
\(I_n(x_n)\subset\mathbb R\) de modo que, para toda sección
\(x=(x_n)_n\),
\[
 I_{n+1}(x_{n+1})\subseteq I_n(x_n),
 \qquad
 \operatorname{diam}I_n(x_n)\longrightarrow0.
\]
Entonces existe un único valor
\begin{equation}
 \operatorname{ev}_{\mathbb R}(x)
 \in\bigcap_{n\geq0}I_n(x_n).
 \label{eq:evaluacion-seccion-continua-rev2}
\end{equation}
Si las asignaciones de intervalos son localmente constantes en cada nivel, la
evaluación \(\operatorname{ev}_{\mathbb R}:\Omega_\infty\to\mathbb R\) es
continua.
\end{theorem}

\begin{proof}
La sobreyectividad permite construir inductivamente prefijos compatibles; la
compacidad de los conjuntos finitos y la propiedad de intersección finita dan
una sección global. El producto \(\prod_nS_n\) es compacto y
\(\Omega_\infty\) es cerrado, pues cada ecuación de compatibilidad define un
cerrado. Para una sección fija, los intervalos son compactos, anidados y de
diámetro tendente a cero; el teorema de intervalos encajados proporciona un
único punto común. Dado \(\varepsilon>0\), se elige un nivel cuyo intervalo
tenga diámetro menor que \(\varepsilon\). Las secciones que comparten ese
prefijo tienen evaluaciones dentro del mismo intervalo, lo que prueba la
continuidad.
\end{proof}

La recta real aparece así como codominio de una contracción. Dos secciones
pueden compartir valor y conservar registros diferentes; la relación
\[
 x\sim_{\mathbb R}y
 \quad\Longleftrightarrow\quad
 \operatorname{ev}_{\mathbb R}(x)=\operatorname{ev}_{\mathbb R}(y)
\]
forma el cociente arquimediano. La genealogía permanece en
\(\Omega_\infty\), mientras la coordenada vive en su imagen real. El problema
del continuo recibe, por tanto, una respuesta constructiva dentro del dominio
cubierto por los sistemas de prefijos declarados: continuidad de la
evaluación y memoria de la sección son propiedades coordinadas, no
descripciones rivales.

\subsection{Cobertura coherente y reconstrucción del cociente arquimediano}

La evaluación anterior asigna un valor a cada sección. Para demostrar que el
espacio visible queda completamente reconstruido se necesita, además, que el
refinamiento no pierda puntos. Sea \(K\subset\mathbb R\) compacto y supóngase
que los intervalos del teorema anterior satisfacen
\begin{align}
 K&=\bigcup_{a\in S_0} I_0(a),
 \label{eq:cobertura-inicial-continuo-rev3}\\
 I_n(a)&=\bigcup_{\substack{b\in S_{n+1}\\r_{n+1,n}(b)=a}}I_{n+1}(b)
 \qquad(a\in S_n).
 \label{eq:cobertura-hijos-continuo-rev3}
\end{align}
Estas identidades definen una \emph{cobertura coherente}: cada compacto
visible se descompone exactamente en los compactos de sus prolongaciones.

\begin{theorem}[Reconstrucción genealógica de un compacto]
\label{thm:cociente-genealogico-compacto-rev3}
Supóngase que:
\begin{enumerate}
\item los conjuntos \(S_n\) son finitos y no vacíos, y los enlaces
      \(r_{n+1,n}\) son sobreyectivos;
\item los intervalos son compactos no vacíos, anidados y uniformemente
      contractivos;
\item se cumplen
      \eqref{eq:cobertura-inicial-continuo-rev3} y
      \eqref{eq:cobertura-hijos-continuo-rev3}.
\end{enumerate}
Entonces \(\operatorname{ev}_{\mathbb R}(\Omega_\infty)=K\), y la aplicación
inducida
\begin{equation}
 \overline{\operatorname{ev}}_{\mathbb R}:
 \Omega_\infty/\!\sim_{\mathbb R}\xrightarrow{\ \cong\ }K
 \label{eq:homeomorfismo-cociente-continuo-rev3}
\end{equation}
es un homeomorfismo. La fibra sobre \(y\in K\) es el conjunto completo de
secciones compatibles que publican ese mismo valor.
\end{theorem}

\begin{proof}
La compacidad y continuidad proceden del
\cref{thm:reconstruccion-continuo-secciones-rev2}. Fijado \(y\in K\), sea
\[
 T_n(y):=\{a\in S_n:y\in I_n(a)\}.
\]
La cobertura coherente hace no vacíos todos los \(T_n(y)\) y conecta cada
estado con una prolongación. El árbol resultante es finitamente ramificado y
posee vértices a toda profundidad; por el lema de la rama infinita admite una
cadena compatible \((x_n)_n\). El punto \(y\) pertenece a todos los
\(I_n(x_n)\), y la contracción fuerza
\(\operatorname{ev}_{\mathbb R}(x)=y\). Por tanto, la evaluación es
sobreyectiva. La aplicación del cociente es continua y biyectiva; como su
dominio es compacto y \(K\) es Hausdorff, es un homeomorfismo.
\end{proof}

El teorema distingue dos afirmaciones. La contracción produce un valor único
para cada genealogía; la cobertura coherente demuestra que todos los valores
del compacto aparecen. Una comprobación a profundidad finita certifica los
niveles examinados, pero no sustituye ninguna de estas hipótesis infinitas.

\begin{proposition}[Atlas genealógico de la recta real]
\label{prop:atlas-genealogico-recta-real-rev3}
Para cada \(m\geq1\), sea \(K_m=[-m,m]\) y supóngase dado un sistema que
satisface el \cref{thm:cociente-genealogico-compacto-rev3}, junto con
inclusiones cerradas
\[
 \jmath_m:\Omega_\infty^{(m)}\hookrightarrow\Omega_\infty^{(m+1)}
\]
que entrelazan las evaluaciones y la inclusión
\(K_m\hookrightarrow K_{m+1}\). Entonces el cociente arquimediano del
colímite estricto de las secciones es
\[
 \varinjlim_m K_m=\bigcup_{m\geq1}[-m,m]=\mathbb R,
\]
con su topología usual, mientras las fibras del colímite conservan las
genealogías de cada valor.
\end{proposition}

\begin{proof}
Cada pieza posee cociente \(K_m\) por
\eqref{eq:homeomorfismo-cociente-continuo-rev3}. Las inclusiones declaradas
hacen compatibles esos cocientes. La topología final del agotamiento por
intervalos compactos coincide con la topología usual de \(\mathbb R\): un
conjunto es abierto exactamente cuando su intersección con cada \(K_m\) es
abierta en \(K_m\).
\end{proof}

\subsection{Descenso algebraico de las operaciones}

Reconstruir el espacio visible no basta para reconstruir su álgebra. Sea
\(\star\) una operación parcial continua sobre \(\Omega_\infty\), con dominio
\(D_\star\) saturado por las fibras de
\(\operatorname{ev}_{\mathbb R}\times\operatorname{ev}_{\mathbb R}\).

\begin{theorem}[Criterio de descenso por fibras]
\label{thm:descenso-operaciones-continuo-rev3}
Existe una operación visible \(\overline\star\) que satisface
\[
 \operatorname{ev}_{\mathbb R}(x\star y)
 =\operatorname{ev}_{\mathbb R}(x)\,\overline\star\,
  \operatorname{ev}_{\mathbb R}(y)
\]
si y sólo si el valor del miembro izquierdo es constante al sustituir
\(x\) o \(y\) por cualquier otro representante de sus respectivas fibras.
Si dos operaciones continuas descendidas coinciden con la suma y el producto
en subconjuntos densos de sus dominios visibles, coinciden con ellas en todo
el dominio.
\end{theorem}

\begin{proof}
La constancia sobre fibras es exactamente la propiedad universal del
cociente: permite definir \(\overline\star\) sin depender del representante,
y es necesaria para que esa definición tenga sentido. La última afirmación
sigue de la continuidad y de la densidad.
\end{proof}

APP conserva el dato requerido para abordar este descenso: las hojas suma y
producto actúan sobre el mismo soporte, pero mantienen cocientes, acarreos y
fibras diferentes. La operación visible puede descender aunque la memoria de
sus representantes siga siendo distinta.

\subsection{Medida proyectiva y conservación de la unidad}

La topología de cilindros organiza proximidad y truncamiento; una teoría de
caminos requiere además pesos. Sea \(\mu_n\) una probabilidad sobre \(S_n\).
La familia es proyectivamente consistente cuando
\begin{equation}
 \mu_n(a)=
 \sum_{\substack{b\in S_{n+1}\\r_{n+1,n}(b)=a}}\mu_{n+1}(b).
 \label{eq:consistencia-medida-cilindrica-rev3}
\end{equation}
La igualdad expresa la conservación local de la unidad: el peso de una
historia truncada es la suma de los pesos de todas sus prolongaciones.

\begin{theorem}[Medida genealógica sobre el límite inverso]
\label{thm:medida-genealogica-limite-rev3}
Toda familia de probabilidades finitas que satisface
\eqref{eq:consistencia-medida-cilindrica-rev3} determina una única
probabilidad de Borel \(\mu_\infty\) sobre \(\Omega_\infty\) tal que
\begin{equation}
 \mu_\infty([a]_n)=\mu_n(a)
 \label{eq:medida-cilindro-rev3}
\end{equation}
para todo cilindro de prefijo \(a\in S_n\).
\end{theorem}

\begin{proof}
Las uniones finitas de cilindros forman un álgebra. Toda una de ellas puede
refinarse hasta una unión disjunta en un mismo nivel, y
\eqref{eq:consistencia-medida-cilindrica-rev3} demuestra que la suma de sus
pesos no depende del nivel elegido. Se obtiene así una premedida. Como los
cilindros son abiertos y cerrados y \(\Omega_\infty\) es compacto, una
sucesión decreciente de elementos del álgebra con intersección vacía alcanza
ya el vacío en algún nivel. La premedida es, por tanto, continua en el vacío
y se extiende de manera única a la \(\sigma\)-álgebra de Borel.
\end{proof}

La evaluación transporta esa medida al compacto visible,
\begin{equation}
 \nu=(\operatorname{ev}_{\mathbb R})_*\mu_\infty,
 \qquad
 \nu(B)=\mu_\infty\bigl(\operatorname{ev}_{\mathbb R}^{-1}(B)\bigr).
 \label{eq:pushforward-medida-visible-rev3}
\end{equation}
La probabilidad visible agrega las genealogías de una misma fibra sin
borrarlas del espacio total.

\begin{proposition}[Desintegración por valores y multiplicidad genealógica]
\label{prop:desintegracion-valores-genealogicos-rev3}
Si \(K\) es un espacio estándar de Borel, existe, para \(\nu\)-casi todo
\(y\in K\), una probabilidad condicionada \(\mu_y\) soportada en
\[
 \mathfrak G_y:=\operatorname{ev}_{\mathbb R}^{-1}(y)
\]
tal que, para toda función integrable \(F\),
\begin{equation}
 \int_{\Omega_\infty}F\,\mathrm d\mu_\infty
 =\int_K
   \left(\int_{\mathfrak G_y}F\,\mathrm d\mu_y\right)
   \mathrm d\nu(y).
 \label{eq:desintegracion-fibras-visibles-rev3}
\end{equation}
Así, la coordenada visible y la distribución de historias que la realizan
son datos tipados distintos: una observable que factoriza por
\(\operatorname{ev}_{\mathbb R}\) sólo lee \(\nu\), mientras un lector de
memoria puede distinguir las medidas \(\mu_y\) dentro de una misma fibra.
\end{proposition}

\begin{proposition}[Condicionamiento compatible de cilindros]
\label{prop:condicionamiento-cilindrico-compatible-rev3}
Sea \(A\subseteq\Omega_\infty\) una unión de cilindros de nivel \(m\), con
\(\mu_\infty(A)>0\). Para todo \(n\geq m\), la distribución condicionada
sobre \(S_n\) se obtiene suprimiendo los estados cuyo truncamiento no
pertenece a \(A\) y renormalizando los pesos restantes. Las distribuciones
resultantes siguen satisfaciendo la consistencia proyectiva
\eqref{eq:consistencia-medida-cilindrica-rev3}.
\end{proposition}

\begin{proof}
Es la regla de Bayes sobre los conjuntos finitos de nivel \(n\). La suma de
los pesos condicionados de las prolongaciones coincide con el peso condicionado del
padre porque la intersección con \(A\) distribuye sobre la unión disjunta de
sus prolongaciones. El condicionamiento elimina futuros incompatibles y
redistribuye la unidad entre los supervivientes; no crea historias nuevas ni
borra del espacio total las que quedaron fuera del suceso leído.
\end{proof}

Esta separación proporciona la base exacta para distinguir preparación,
lectura, condicionamiento e historia en la realización cuántica.

\begin{statusbox}[title={Fuerza matemática y función arquitectónica}]
El rango \(n-1\), el teorema fundamental discreto, la identidad cocíclica del
acarreo, los teoremas de reconstrucción y la medida proyectiva son resultados
exactos en los dominios declarados. Su aplicación HMT utiliza como estructura de partida los estados,
restricciones, intervalos y transportes construidos por APP--TRIT--TPK. El
complejo discreto de medida es una formalización organizadora: no identifica
el cociclo categorial de grado dos con una celda geométrica hasta que se
proporcione el mapa correspondiente. Esta separación permite que el lenguaje
común gobierne después la holonomía, la suma sobre caminos y la transición de
número a partícula sin anticipar sus realizaciones físicas.
\end{statusbox}
